\documentclass[11pt,letterpaper]{article}

\usepackage[margin=1in]{geometry}
\usepackage{termcal}
\usepackage{enumitem}
\usepackage[colorlinks=true, allcolors=blue]{hyperref}
\usepackage{color}
\usepackage{multirow}
\usepackage{multicol}

\newcommand{\todo}[1]{\textcolor{red}{TODO: #1}}

\title{18-874: Spacecraft Design-Build-Fly Lab}
\author{Spring 2026}
\date{}

\begin{document}

\maketitle

\section*{Course Description}

Spacecraft design is a truly interdisciplinary subject that draws from every branch of engineering. This capstone design class brings together the material from prior classes in a way that emphasizes the interactions between disciplines and demonstrates how some of the more theoretical topics are synthesized in the practical design of a spacecraft. The class will design, build, and test a small satellite that addresses objectives and requirements posed at the beginning of the course sequence. Students will work in subsystem teams, each focusing on some aspect of the spacecraft, but exposed to many different disciplines and challenges. Practical, hands-on engineering skills will be emphasized, along with fabrication and testing of physical hardware and the creation of thorough documentation.

\subsection*{Overview for Fall 2026:} In this offering of the course the team
will be developing and implementing features for Argus-3, based on the design
of Argus-2, which is currently in the final process of launch logistics for a future launch in 2027/2028 timeframe.  
%
These Argus-3 features will be designed, developed and integrated into a future-looking variant of the Argus flight
bus that will fly shortly after the Argus-2 bus.

\section*{Instructors}

\begin{center}
\begin{tabular}{l l}
	Prof. Brandon Lucia & \textbf{Email:} \href{mailto:blucia@andrew.cmu.edu}{blucia@andrew.cmu.edu}
	\\
	System Scientist \& Labmaster: Neil Khera & \textbf{Email:} \href{nkhera@andrew.cmu.edu}{nkhera@andrew.cmu.edu} \\
	TA: Sun A Cho & \textbf{Email:} \href{sunac@andrew.cmu.edu}{sunac@andrew.cmu.edu} \\
	TA: Om Anavekar & \textbf{Email:} \href{oanaveka@andrew.cmu.edu}{oanaveka@andrew.cmu.edu} \\
	TA: Veronica Muriga & \textbf{Email:} \href{vmuriga@andrew.cmu.edu}{vmuriga@andrew.cmu.edu} 
\end{tabular}
\end{center}

\section*{Learning Objectives}

The goal of this course is to give students hands-on experience designing and
building small spacecraft subsystems and integrating them into a CubeSat.
Throughout this course, students will:
\begin{enumerate}
	\item Understand how the design and integration of a system whose performance depends on the success of many interacting subsystems.
	\item Work within a small team to fabricate, and test hardware and software through rapid design iteration.
	\item Coordinate with other teams to integrate subsystems into a complete spacecraft.
	\item Gain exposure to the complete life cycle of a small satellite mission.
	\item Develop skills at project management and systems engineering
	\item Develop social skills and planning skills to work with a large group and coordinate many required tasks to launch-ready a satellite.
\end{enumerate}


\section*{Logistics}

The course will involve designing and building hardware in small teams. Class time will be used primarily for weekly team meetings and consulting time to meet with the instructors.

\begin{itemize}
	\item Class sessions: MW 3:30-4:50
	\item All-hands meetings will be held every week on Wednesdays, followed by consulting hours.
	\item Mondays will be in-lab working sessions unless there is a lecture scheduled, in which case lecture will be in PH A18A
	\item Students are expected to be in lab during in-lab working sessions and outside of schedule course hours.
	\item Slack will be used for coordination between teams and instructors. All students will be added to the ``SpacecraftDesignBuildFlyLab'' slack channel.
	\item GitHub will be used to manage project files for all teams and teams will be required to use good git hygiene and software engineering practices.
\end{itemize}

\section*{Assignments and Exams}

There will be no exams in this course. Evaluation will be based on participation, contribution to design and fabrication work, final documentation, and weekly project management metrics from each team.  Grade details are broken down below.  Use of project management tools is mandatory and your grade will be derived from specific documentation of technical contributions.

You will keep an individual weekly lab notebook.  This will document your progress and activities for the week, every week.  Your team will additionally keep an all-hands presentation log that gets updated on a weekly cadence as well.

\section*{Grading}

Grading will be based on:
\begin{itemize}
	\item 20\% Participation and attendance of team meetings
	\item 30\% Individual technical contributions quantified by git commit history / number and significance of PRs and peer surveys
	\item 30\% Completeness and quality of team documentation
	\item 20\% Final presentation and report on team progress for semester
\end{itemize}


\section*{Learning Resources}

There is no textbook required for this course.

\section*{Course Policies}

\textbf{Attendance:} This is a team-based course. In order to coordinate work among teams, participation in weekly meetings is required. If you are unable to be present at a meeting, you must notify the instructors and ensure that your teammates are prepared to present your work.

\medskip
\noindent
\textbf{Accommodations for Students with Disabilities:} If you have a disability and are registered with the Office of Disability Resources, I encourage you to use their online system to notify me of your accommodations and discuss your needs with me as early in the semester as possible. I will work with you to ensure that accommodations are provided as appropriate. If you suspect that you may have a disability and would benefit from accommodations but are not yet registered with the Office of Disability Resources, I encourage you to contact them at \href{mailto:access@andrew.cmu.edu}{access@andrew.cmu.edu}.

\medskip
\noindent
\textbf{Statement of Support for Students' Health \& Well-Being:} Take care of yourself. Do your best to maintain a healthy lifestyle this semester by eating well, exercising, avoiding drugs and alcohol, getting enough sleep, and taking some time to relax. This will help you achieve your goals and cope with stress.

\medskip
\noindent
If you or anyone you know experiences any academic stress, difficult life events, or feelings like anxiety or depression, we strongly encourage you to seek support. Counseling and Psychological Services (CaPS) is here to help: call 412-268-2922 and visit \href{http://www.cmu.edu/counseling}{http://www.cmu.edu/counseling}. Consider reaching out to a friend, faculty, or family member you trust for help getting connected to the support that can help.

\medskip
\noindent
\textit{If you or someone you know is feeling suicidal or in danger of self-harm, call someone immediately, day or night:}

\textit{CaPS: 412-268-2922}

\textit{Re:solve Crisis Network: 888-796-8226}

\medskip
\noindent
\textit{If the situation is life threatening, call the police:}

\textit{On campus: CMU Police: 412-268-2323}

\textit{Off campus: 911}


\section*{Tentative Fall 2026 Schedule}

\begin{tabular}{c|c|c}
	Week & Dates & Topics \\
	\hline
	\multirow{2}{*}{1} & Aug 24 & Course Overview \& Logistics \\
	 & Aug 26 & Mission PDR Review and Team Assignments \\
	\hline
	\multirow{2}{*}{2} & Aug 31 & Project Management \\
	 & Sep 2 & All-hands Meeting\\
	\hline
	\multirow{2}{*}{3}  & Sep 7 & \textcolor{red}{No Class (Labor Day)} \\
	 & Sep 9 & All-Hands Meeting \\
	\hline
	\multirow{2}{*}{4} & Sep 14 & Lecture Orbital Edge Computing \\
	 & Sep 16 & All-Hands Meeting \\
	\hline
	\multirow{2}{*}{5}  & Sep 21 & \textcolor{red}{No Class (Instructor Travel)} \\
	 & Sep 23 & All-Hands Meeting  \\
	\hline
	\multirow{2}{*}{6}  & Sep 28 & Lecture Efficient's Dataflow Architecture  \\
	 & Sep 30 & All-Hands Meeting \\
	\hline
	\multirow{2}{*}{7}  & Oct 5 &  Case Study: Tartan Artibeus \\
	 & Oct 7 & All-Hands Meeting \\
	\hline
	\multirow{2}{*}{8}  & Oct 12 & \textcolor{red}{No Class (Spring Break)} \\
	 & Oct 14 & \textcolor{red}{No Class (Spring Break)} \\
	\hline
	\multirow{2}{*}{9}  & Oct 19 & Guest Lecture: Paulo Fisch \\
	 & Oct 21 & All-Hands Meeting \\
	\hline
	\multirow{2}{*}{10}  & Oct 26 & Case Study: PROMETHEUS-1 (Neil Khera) \\
	 & Oct 28 & All-Hands Meeting \\
	 \hline
	\multirow{2}{*}{11}  & Nov 2  & Lab Work Day \\
	 & Nov 4 & All-Hands Meeting \\
	 \hline
	\multirow{2}{*}{12}  & Nov 9 & Guest Lecture or Lab Day \\
	 & Nov 11 & All-Hands Meeting \\
	 \hline
	\multirow{2}{*}{13}  & Nov 16 &\textcolor{red}{No Class (Instructor Travel)} \\
	 & Nov 18 & All-Hands Meeting (Prof. Lucia out) \\
	 \hline
	\multirow{2}{*}{14}  & Nov 23 & Final presentation feedback and final lab day \\
	 & Nov 25 & \textcolor{red}{No Class (Thanksgiving week)} \\
	 \hline
	\multirow{2}{*}{14}  & Nov 30  & Final Presentations / Design Review \\
	 & Dec 2 & Final Presentations / Design Review \\
\end{tabular}


\end{document}
